\documentclass[a4paper]{report}
% Shared preamble for the English probability and statistics documents.
% Based on the original course preamble. Compile with pdfLaTeX.
\usepackage[top=1.5cm,bottom=2cm,left=2cm,right=2cm]{geometry}
\usepackage[T1]{fontenc}
\usepackage{lmodern}
\usepackage[utf8]{inputenc}
\usepackage[english]{babel}

\usepackage{amsmath}
\usepackage{amssymb}
\usepackage{amsthm}
\usepackage{mathrsfs}
\usepackage{mathtools}
\usepackage{etoolbox}
\usepackage{tikz}
\usetikzlibrary{shapes.multipart,arrows.meta}
\usepackage{xcolor}
\usepackage{enumerate}
\usepackage[shortlabels]{enumitem}
\usepackage{hyperref}
\usepackage{version}
\usepackage{bbold}
\usepackage{fontawesome5}
\usepackage{fancyhdr}
\lfoot{}
\usepackage{natbib}
\usepackage{graphicx}
\usepackage{pdfpages}
\usepackage{authblk}
    
\usepackage{blindtext}
\usepackage{multicol}
\usepackage{collcell}
\usepackage{wrapfig}
\usepackage{blkarray}
\usepackage{cancel}
\usepackage{appendix}
\usepackage{float}
\makeatletter
\input{tex/assets/slashbox.sty}
\makeatother
\hypersetup{unicode,colorlinks,citecolor=red,filecolor=black,
  linkcolor=red,urlcolor=cyan}
\setlength{\columnsep}{1cm}
\pagenumbering{gobble}

\def\numtext#1{%
  \ifcase#1\or one\or two\or three\or four\or five\or six\or seven\or eight\or nine\or ten\or
  eleven\or twelve\or thirteen\or fourteen\or fifteen\or sixteen\or seventeen\or eighteen\or nineteen\or twenty\or Lots\fi}

\newcommand{\Byte}[1]{%
  \def\nodecontents{}%
  \foreach[count=\l] \k in {#1} {\xappto{\nodecontents}{\noexpand\nodepart{\numtext{\l}}\k}}%
  \node[rectangle split,rectangle split parts=8,rectangle split horizontal,draw] at (2,2){\nodecontents};}

\newcommand{\D}{{\mathrm{d}}}
\newcommand{\cE}{{\mathcal{E}}}
\newcommand{\cG}{{\mathcal{G}}}
\newcommand{\cX}{{\mathcal{X}}}
\newcommand{\cY}{{\mathcal{Y}}}
\newcommand{\cB}{{\mathcal{B}}}
\newcommand{\cF}{{\mathcal{F}}}
\newcommand{\cO}{{\mathcal{O}}}
\newcommand{\cS}{{\mathcal{S}}}
\newcommand{\cT}{{\mathcal{T}}}
\newcommand{\given}{\, \big| \,}
\newcommand{\cL}{{\mathscr{L}}}
\newcommand{\cD}{{\mathcal{D}}}
\newcommand{\cA}{{\mathscr{A}}}
\newcommand{\cP}{{\mathscr{P}}}
\newcommand{\cR}{{\mathscr{R}}}
\newcommand{\A}{{\mathcal A}}
\newcommand{\one}{\mathbb{1}}
\newcommand{\mathbbm}[1]{\mathbb{#1}}

\providecommand{\C}{}
\renewcommand{\C}{\mathbb C}
\newcommand{\R}{\mathbb R}
\newcommand{\Q}{\mathbb Q}
\newcommand{\N}{\mathbb N}
\newcommand{\T}{\mathbb T}
\newcommand{\F}{\mathbb F}
\renewcommand{\P}{\mathbb P}
\newcommand{\Z}{\mathbb Z}
\newcommand{\E}{\mathbb {E}}
\newcommand{\Pbb}[1]{\P\left( #1\right)}
\newcommand{\Ebb}[1]{\E\left[ #1\right]}

\newcommand{\lb}{\left (}
\newcommand{\rb}{\right )}
\newcommand{\lbb}{\left [}
\newcommand{\rbb}{\right ]}
\newcommand{\labs}{\left |}
\newcommand{\rabs}{\right |}
\newcommand{\lbrb}[1]{\lb #1 \rb}
\newcommand{\lbbrbb}[1]{\lbb#1\rbb}
\newcommand{\labsrabs}[1]{\labs#1\rabs}
\newcommand{\lbbrb}[1]{\lbb#1\rb}
\newcommand{\lbrbb}[1]{\lb#1\rbb}
\newcommand{\langlerangle}[1]{\langle#1\rangle}
\newcommand{\lbcurly}{\{}
\newcommand{\rbcurly}{\}}
\newcommand{\curly}[1]{\lbcurly#1\rbcurly}
\newcommand{\abs}{\labsrabs}
\newcommand{\ind}[1]{1_{#1}}
\newcommand{\indep}{\perp \!\!\! \perp}
\renewcommand{\d}{\, \mathrm{d}}

\newtheorem{theorem}{Theorem}
\newtheorem{problemCnt}{Lemma}
\theoremstyle{definition}
\newtheorem{problem}[problemCnt]{Problem}
\newtheorem{problembook}[problemCnt]{\faBook \,  Problem}
\newtheorem{problempencil}[problemCnt]{\faPencil* \,  Problem}
\newtheorem{problemstar}[problemCnt]{\faStar\,  Problem}
\newtheorem{problemstarpencil}[problemCnt]{\faStar \, \faPencil* \,   Problem}
\newtheorem{probleminteresting}[problemCnt]{\faSmile[regular] \,   Problem}
\newtheorem{solution}{Solution}

\newcommand{\fwcell}[1]{\makebox[\arraystretch\normalbaselineskip]{$#1$}}
\newenvironment{enumeratei}{\begin{enumerate}[(i)]}{\end{enumerate}}
\newenvironment{enumeratea}{\begin{enumerate}[(a)]}{\end{enumerate}}
\newenvironment{enumeraten}{\begin{enumerate}[\upshape 1.]}{\end{enumerate}}

% Compatibility helpers needed by a few files in the combined archive.
\AtBeginDocument{\noindent{\footnotesize English translation prepared with ChatGPT (OpenAI, GPT-6), 26 October 2026, for translation of the Bulgarian text into English and LaTeX formatting.
The Bulgarian originals remain the reference; translation may contain errors.}\par\smallskip}
\newcommand{\MissingGraphic}[1]{%
  \begin{center}\fbox{\parbox{0.82\linewidth}{\centering Missing source image: \texttt{#1}}}\end{center}}



\setlength{\emergencystretch}{3em}
\allowdisplaybreaks
\begin{document}



\vspace{1ex}
\noindent Probability and Statistics, 19.01.2025 \hfill Test 2

\hrulefill

\begin{it}
The exact grading formula depends on the results. As an approximation, use 2 plus the points earned. Time allowed: 210 minutes. Good luck.
\end{it}
\vspace{1ex}


\begin{problem} (1.5 points)
Let $\xi,\eta$ be independent random variables,
$\xi\sim Exp(4)$ and $\eta\sim U(-1,1)$.
\begin{enumerate}
    \item (0.5 points) Find
    $\Ebb{\xi \eta}, \Ebb{\xi|
    \eta=0}$, $\mathrm{Corr}(\xi,\eta)$ and $P(\xi\leq0.25)$.
    \item (0.5 points) Find the densities of $\xi^2+3$ and $\xi/
\eta$.
\item (0.5 points)
Let $Y_1,Y_2,Y_3,Y_4\sim Exp(1)$ be independent. Let $Y_{(i)}$ be the $i$th smallest.
Find the expectation and variance of $Y_{(3)}$.
\end{enumerate}
\end{problem}


\begin{problem}
(1 point)
Two customers arrive simultaneously at a service system
with two queues. If $T$
is the service time,
\[
T \sim
\begin{cases}
Exp(1), &\text{(slow queue) with probability }0.3,\\
Exp(3), &\text{(fast queue) with probability }0.7.
\end{cases}
\]
\begin{enumerate}
    \item (0.5 points) Let $T_i$ be customer $i$'s service time. Find the distribution of $T_1+T_2$.
    \item (0.25 points) Find $E(T_1+T_2)$.
    \item (0.25 points) For large $n$, estimate the probability that the sum of $n$ independent such service times
    exceeds $5n$.
\end{enumerate}
\end{problem}


\begin{problem}
(1 point)
Let $X_1,X_2,\dots$ be independent identically distributed random variables
with density
\[
f_X(x) =
\begin{cases}
2x, &\text{if }0<x<1,\\
0, &\text{otherwise}.
\end{cases}
\]
\begin{enumerate}
    \item (0.5 points) Approximately how many observations $n$ are needed for a 97\% probability
    that the sample mean $\overline X_n:=(X_1+\dots+X_n)/n$ lies between 0.5 and 0.7?
    \item (0.5 points) Approximately what value do you expect for $\max_{1\leq k\leq n}X_k$ when $n=10\,000$?
\end{enumerate}
\end{problem}


\begin{problem}
(0.75 points)
Let $X,Y$ be independent, with $X\sim\mathcal N(2,4)$ and
$Y \sim \mathcal{N}(2, 9)$.
\begin{enumerate}
    \item (0.25 points) Find their joint density and $P(X<Y)$.
    \item (0.5 points) Calculate $E(|X-Y|)$ and $|E(X)-E(Y)|$.
    Can you determine which is larger without computing their exact values?
\end{enumerate}
\end{problem}


\begin{problem}
(1 point)
Let $ABC$ be a right triangle with legs $AC=BC=1$.
Through a uniformly chosen point $P$ on $AC$, draw a line parallel to hypotenuse $AB$.
Let $Q,R$ be its intersections with the legs.
\begin{enumerate}
    \item (0.25 points) Find the expected length of $QR$.
    \item (0.5 points) Find the probability that two independently chosen random points inside $ABC$ lie on opposite sides of $QR$.
    \item (0.25 points) Let $A,B$ be randomly chosen on
    a circle with centre $O$ and radius 1. What is the expected area of $AOB$? What is the expectation of $AO\cdot AB\cdot\sin(OAB)/2$? Explain why the two quantities are or are not equal.
\end{enumerate}
\end{problem}
\pagebreak
\begin{problem}
\textit{(Bonus: 1 point)}

In classical (often called frequentist) statistics, to estimate an unknown parameter \(\lambda\),
we usually use methods such as maximum likelihood or the method of moments,
without modelling \(\lambda\) itself as random.
In the \emph{Bayesian approach}, we instead assume a \emph{prior}
distribution for \(\lambda\), such as \(\Gamma(\alpha,\beta)\), and combine observed data with this prior knowledge
to obtain its \emph{posterior} distribution.
We can then generate or predict new observations using a distribution that accounts for data already observed.

\bigskip
Models in which the posterior belongs to the same family as the prior, with updated parameters, are computationally convenient: the data
have taught us parameter values that better describe reality. One such case is the exponential model
$Exp(\lambda)$ with a gamma prior for the unknown $\lambda$: its posterior is also gamma.

Let \(X_1,\dots,X_n\) be independent observations from an exponential distribution,
\[
X\mid\lambda \sim \mathrm{Exp}(\lambda),
\]
with unknown parameter \(\lambda>0\). In the Bayesian model, assume the prior \(\lambda\sim\Gamma(\alpha,\beta)\), with density
\[
f_\lambda(x) \;=\; \frac{\beta^\alpha}{\Gamma(\alpha)}\,x^{\alpha-1}\,\exp(-\beta x),
\quad x>0.
\]

\begin{enumerate}
    \item (0.25 points) Prove that, after observing \(X_1,\dots,X_n\), the posterior distribution of \(\lambda\) is
    \[\lambda\mid X_1,\ldots,X_n\sim\mathrm{\Gamma}\Bigl(\alpha + n,\;\beta + \sum_{i=1}^n X_i\Bigr).\]

    \item (0.25 points) Find the predictive distribution of a new observation \(X_{n+1}\)
    (i.e. the conditional probability \(\mathbb P(X_{n+1}\in A|X_1,\dots,X_n)\))
    and show that in certain cases it can be expressed using a \textit{Pareto} random variable $Y$:
    \[
    F_Y(x) = 1 - \lbrb{\frac{a}{x}}^\alpha,
    \quad
    x > a,
    \]
    for a parameter $a$. This is
    a classical example of a heavy-tailed distribution.
    (\textit{You may integrate with respect to \(x\) the distribution
    \(\mathbb{P}(X_{n+1}\in A|\lambda=x)\times f_{\lambda\mid X_1,\dots,X_n}(x|X_1,\ldots,X_n)\).})

    \item (0.5 points)
    Find \(\mathrm E(\lambda|X_1,\dots,X_n)\) and \(\mathrm{Var}(\lambda|X_1,\dots,X_n)\).
    Recall that the conditional variance of \(Y\) given \(X\) is defined as
    \[
       D(Y|X) \;=\; E\bigl[(Y - E(Y|X))^2 \,\big|\;X\bigr].
    \]
    Use this definition also to prove that
    \[
       D(Y) \;=\; E\bigl[D(Y|X)\bigr]\;+\; D\bigl(E[Y|X]\bigr).
    \]

    \item \textit{(Bonus to the bonus)}
    How would you interpret the difference between a classical confidence interval and
    a Bayesian highest-posterior-density interval for \(\lambda\) in this example?
\end{enumerate}
\end{problem}
















\end{document}
